\documentclass[12pt,a4paper]{amsart}
\usepackage{amsmath,amssymb,amsthm}
\usepackage{mathtools}
\usepackage[utf8]{inputenc}
\usepackage[T1]{fontenc}
\usepackage[ngerman]{babel}
\usepackage{hyperref}
\usepackage{enumitem,booktabs,array}
\usepackage{geometry}
\geometry{margin=2.5cm}

% -------------------------------------------------------
% Theorem-Umgebungen (Deutsch)
% -------------------------------------------------------
\newtheorem{theorem}{Satz}[section]
\newtheorem{lemma}[theorem]{Lemma}
\newtheorem{corollary}[theorem]{Korollar}
\newtheorem{proposition}[theorem]{Proposition}
\theoremstyle{definition}
\newtheorem{definition}[theorem]{Definition}
\newtheorem{example}[theorem]{Beispiel}
\newtheorem{remark}[theorem]{Bemerkung}
\newtheorem{conjecture}[theorem]{Vermutung}

% -------------------------------------------------------
% Eigene Befehle
% -------------------------------------------------------
\newcommand{\N}{\mathbb{N}}
\newcommand{\Z}{\mathbb{Z}}
\newcommand{\Q}{\mathbb{Q}}
\newcommand{\R}{\mathbb{R}}
\newcommand{\C}{\mathbb{C}}
\DeclareMathOperator{\lcm}{kgV}
\DeclareMathOperator{\ggT}{ggT}
\DeclareMathOperator{\ord}{ord}

% -------------------------------------------------------
\begin{document}
% -------------------------------------------------------

\title{Die Erd\H{o}s--Straus-Vermutung:\\
       Ägyptische Brüche, Diophantische Gleichungen\\
       und partielle Resultate}

\author{Michael Fuhrmann}
\address{Specialist Mathematics Project, Build 133}
\email{specialist-maths@localhost}
\date{\today}

\begin{abstract}
Die Erd\H{o}s--Straus-Vermutung, formuliert von Paul Erd\H{o}s und Ernst
Gabor Straus im Jahre 1948, besagt, dass es für jede ganze Zahl $n \geq 2$
positive ganze Zahlen $a, b, c$ gibt mit
\[
  \frac{4}{n} = \frac{1}{a} + \frac{1}{b} + \frac{1}{c}.
\]
Dies ist eine Aussage über ägyptische Brüche (Darstellung rationaler Zahlen
als Summen von Stammbrüchen) und eines der bekanntesten offenen Probleme der
elementaren Zahlentheorie.  Die Vermutung wurde rechnerisch für alle
$n \leq 10^{14}$ verifiziert und ist für fast alle $n$ im Sinne der
Dichte bekannt (Schinzel 1956; Vaughan 1970).  Diese Arbeit stellt den
historischen Hintergrund, alle wichtigen konstruktiven Ansätze, modulare
Reduktionen, Dichteergebnisse, die Verbindung zur Diophantischen Gleichung
$4abc = n(ab+ac+bc)$ sowie die rechnerische Verifikation für alle
Primzahlen bis $10^9$ vor.  Die Vermutung ist offen.
\end{abstract}

\maketitle
\tableofcontents

% =========================================================
\section{Einleitung}
% =========================================================

Eine der ältesten Traditionen der Zahlentheorie ist die Darstellung
rationaler Zahlen als Summen von Stammbrüchen $1/k$, $k \in \N$.
Solche Darstellungen wurden bereits von den alten Ägyptern verwendet
und erscheinen zum Beispiel im Rhind-Papyrus (ca.\ 1650~v.\,Chr.).
Ein \emph{Stammbruch} (oder \emph{ägyptischer Bruch}) ist ein Bruch
der Form $1/k$ mit $k$ einer positiven ganzen Zahl.

Im Jahr 1948 stellten Paul Erd\H{o}s und Ernst Gabor Straus die folgende
Frage: Gilt es für jeden Bruch $4/n$, ihn als Summe von genau drei
Stammbrüchen zu schreiben?  Diese scheinbar einfache Frage ist seit
mehr als siebzig Jahren unbeantwortet.

\begin{conjecture}[Erd\H{o}s--Straus, 1948]\label{verm:ES}
Für jede ganze Zahl $n \geq 2$ existieren positive ganze Zahlen
$a, b, c$ derart, dass
\[
  \frac{4}{n} = \frac{1}{a} + \frac{1}{b} + \frac{1}{c}.
\]
\end{conjecture}

Die Schwierigkeit liegt in der Universalität der Aussage: Wir benötigen
eine Darstellung für \emph{jedes} $n$, nicht nur für eine große Klasse.
Obwohl die Vermutung rechnerisch bis $n = 10^{14}$ ohne eine einzige
Ausnahme überprüft wurde, bleibt ein Beweis fern.

Das Problem steht exemplarisch für die Gattung Diophantischer Probleme,
die leicht zu formulieren, für einzelne Werte leicht zu überprüfen,
jedoch außerordentlich schwer vollständig zu lösen sind.
Es gehört neben der Goldbach-Vermutung und der Collatz-Vermutung zu den
bekanntesten offenen Problemen, die auch für Nicht-Spezialisten zugänglich sind.

\subsection*{Gliederung dieser Arbeit}

Abschnitt~2 behandelt die Geschichte der ägyptischen Brüche.
Abschnitt~3 gibt die formale Aussage und äquivalente Umformulierungen.
Abschnitt~4 entwickelt die drei wichtigsten konstruktiven Ansätze.
Abschnitt~5 behandelt modulare Reduktionen.
Abschnitt~6 deckt Dichte- und analytische Resultate ab.
Abschnitt~7 fasst die rechnerische Verifikation zusammen.
Abschnitt~8 listet verwandte offene Probleme auf.
Abschnitt~9 enthält das Literaturverzeichnis.

% =========================================================
\section{Ägyptische Brüche: Geschichte und Grundlagen}
% =========================================================

\subsection{Antike Ursprünge}

Die Ägypter schrieben keine Brüche $p/q$ mit $p > 1$; stattdessen
stellten sie jede rationale Zahl als endliche Summe von
\emph{verschiedenen} Stammbrüchen dar.  Der Rhind-Papyrus
(ca.\ 1650~v.\,Chr.) enthält eine Tabelle von Zerlegungen von $2/n$
für ungerade $n$ von $3$ bis $101$.  Zum Beispiel:
\[
  \frac{2}{5} = \frac{1}{3} + \frac{1}{15}, \quad
  \frac{2}{7} = \frac{1}{4} + \frac{1}{28}, \quad
  \frac{2}{13} = \frac{1}{8} + \frac{1}{52} + \frac{1}{104}.
\]
Diese Tradition wurde von Fibonacci (Leonardo von Pisa) fortgeführt,
der in seinem \emph{Liber Abaci} (1202) einen Greedy-Algorithmus --
heute \emph{Fibonacci-Sylvester-Algorithmus} genannt -- zur Entwicklung
beliebiger rationaler Zahlen in Stammbrüche beschrieb.

\subsection{Der Fibonacci-Sylvester-Greedy-Algorithmus}

\begin{theorem}[Fibonacci-Sylvester-Algorithmus]\label{satz:sylvester}
Jede positive rationale Zahl $r = p/q$ mit $0 < p/q \leq 1$ besitzt
eine endliche Darstellung als Summe verschiedener Stammbrüche.
\end{theorem}

Der Algorithmus iteriert wie folgt: In jedem Schritt ersetze $p/q$ durch
$p/q - 1/\lceil q/p \rceil$.  Entscheidend ist, dass dieser Prozess
terminiert, weil der Zähler bei jedem Schritt strikt kleiner wird.

\begin{example}
Wir entwickeln $4/7$ mit dem Greedy-Algorithmus:
\[
  \frac{4}{7}: \quad \left\lceil \frac{7}{4} \right\rceil = 2,
  \quad \frac{4}{7} - \frac{1}{2} = \frac{1}{14}.
\]
Also $4/7 = 1/2 + 1/14$.  Es werden nur zwei Stammbrüche benötigt;
Vermutung~\ref{verm:ES} fordert höchstens drei, also ist $4/7$ erledigt.
\end{example}

\begin{example}
Entwicklung von $4/11$:
\[
  \left\lceil \frac{11}{4} \right\rceil = 3,\quad
  \frac{4}{11} - \frac{1}{3} = \frac{1}{33}.
\]
Also $4/11 = 1/3 + 1/33$.  Auch hier nur zwei Brüche.
\end{example}

\subsection{Grundlegende Definitionen}

\begin{definition}[Ägyptische Bruchdarstellung]
Eine \emph{$m$-gliedrige ägyptische Bruchdarstellung} einer positiven
rationalen Zahl $r$ ist eine Multimenge $\{a_1, \ldots, a_m\} \subset \N$ mit
\[
  r = \sum_{i=1}^{m} \frac{1}{a_i}.
\]
Sind zusätzlich alle $a_i$ verschieden, spricht man von einer
\emph{eigentlichen} ägyptischen Bruchdarstellung.
\end{definition}

\begin{remark}
Die Erd\H{o}s--Straus-Vermutung verlangt \emph{nicht}, dass die drei
Nenner $a, b, c$ verschieden sind.  Lässt man Wiederholungen zu,
ist das Problem etwas leichter, doch die Vermutung ist in beiden
Formulierungen von Interesse.  Die Standardformulierung erlaubt $a = b = c$
(was hier aus arithmetischen Gründen unmöglich ist: $4/n = 3/a$ würde
$n \mid 3$ erfordern, was separat behandelt wird).
\end{remark}

% =========================================================
\section{Die Vermutung: Formale Aussage und Reduktionen}
% =========================================================

\subsection{Formale Aussage}

Wir wiederholen die Vermutung in präziser Form.

\begin{conjecture}[Erd\H{o}s--Straus]\label{verm:ES2}
Für jede ganze Zahl $n \geq 2$ existieren $a, b, c \in \N$ mit
\begin{equation}\label{gl:ES}
  \frac{4}{n} = \frac{1}{a} + \frac{1}{b} + \frac{1}{c}.
\end{equation}
\end{conjecture}

\subsection{Reduktion auf Primzahlen}

\begin{proposition}\label{prop:primreduktion}
Es genügt, Vermutung~\ref{verm:ES2} für alle Primzahlen $p \geq 5$
zu beweisen.
\end{proposition}

\begin{proof}
Die kleinen Fälle werden direkt verifiziert:
\begin{align*}
  \frac{4}{2} &= 1 + \frac{1}{2} + \frac{1}{2}, \\
  \frac{4}{3} &= 1 + \frac{1}{4} + \frac{1}{12}, \\
  \frac{4}{4} &= \frac{1}{2} + \frac{1}{4} + \frac{1}{4}.
\end{align*}
Sei nun $n \geq 5$ zusammengesetzt, $n = p \cdot q$ mit Primteiler $p$.
Gilt $4/p = 1/a + 1/b + 1/c$, so folgt
\[
  \frac{4}{n} = \frac{4}{pq} = \frac{1}{qa} + \frac{1}{qb} + \frac{1}{qc},
\]
eine gültige Darstellung für $n$.
\end{proof}

\begin{corollary}
Die Erd\H{o}s--Straus-Vermutung ist äquivalent zu ihrer Einschränkung auf
Primzahlen.
\end{corollary}

\subsection{Die Diophantische Gleichungsform}

Multipliziert man Gleichung~\eqref{gl:ES} mit $nabc$, erhält man die
äquivalente Diophantische Gleichung:

\begin{proposition}[Diophantische Formulierung]\label{prop:dioph}
Gleichung~\eqref{gl:ES} gilt genau dann, wenn
\begin{equation}\label{gl:dioph}
  4abc = n(bc + ac + ab).
\end{equation}
\end{proposition}

Diese Form eignet sich zur Lösungszählung und zur Anwendung der Theorie
quadratischer Formen und multiplikativer Funktionen.

\begin{example}
Für $n = 5$: $4/5 = 1/2 + 1/4 + 1/20$.  Probe:
$4 \cdot 2 \cdot 4 \cdot 20 = 640$ und $5(4 \cdot 20 + 2 \cdot 20 + 2 \cdot 4)
= 5 \cdot 128 = 640$. \checkmark
\end{example}

% =========================================================
\section{Äquivalente Formulierungen}
% =========================================================

\subsection{Zweigliederige Reduktionen}

Manchmal zerlegt sich $4/n$ in eine Summe aus nur zwei Stammbrüchen,
was für die Vermutung mehr als ausreichend ist.

\begin{proposition}[Zweigliedriger Fall]\label{prop:zweigliedrig}
Falls $a, b \in \N$ mit $4/n = 1/a + 1/b$ existieren, gilt die Vermutung
trivialerweise für $n$ (man teile einen der Terme weiter auf,
z.\,B.\ $1/a = 1/(2a) + 1/(2a)$).
\end{proposition}

Die Gleichung $4/n = 1/a + 1/b$ ist äquivalent zu
$4ab = n(a+b)$, also $(4a - n)(4b - n) = n^2$.
Man braucht eine Faktorisierung $n^2 = d_1 d_2$ mit
$d_1 \equiv d_2 \equiv n \pmod{4}$.

\subsection{Umformulierung über Kehrwerte arithmetischer Progressionen}

Eine weitere äquivalente Form: Vermutung~\ref{verm:ES2} gilt für $n$ genau
dann, wenn die Gleichung
\[
  4x_1 x_2 x_3 - n(x_1 x_2 + x_1 x_3 + x_2 x_3) = 0
\]
eine positive ganzzahlige Lösung $(x_1, x_2, x_3)$ besitzt.

\subsection{Multiplikativ-zahlentheoretische Sichtweise}

Für festes $n$ definiere
\[
  N_3(n) = \#\{(a,b,c) \in \N^3 : a \leq b \leq c,\;
            \tfrac{4}{n} = \tfrac{1}{a} + \tfrac{1}{b} + \tfrac{1}{c}\}.
\]
Die Vermutung behauptet $N_3(n) \geq 1$ für alle $n \geq 2$.

Heuristische Zählargumente legen nahe, dass
$N_3(n) \approx C n (\log n)^2$ für eine Konstante $C > 0$ gilt,
was es extrem unwahrscheinlich macht, dass ein $n$ eine Ausnahme ist.

% =========================================================
\section{Konstruktive Ansätze}
% =========================================================

Wir stellen drei grundlegende konstruktive Strategien vor, von denen jede
eine große (aber nicht vollständige) Restklasse von $n$ abdeckt.

\subsection{Typ I: Teilbarkeitskonstruktionen}

\begin{proposition}[Typ I]\label{prop:typI}
Ist $4 \mid n$, so schreibe $n = 4m$. Dann:
\[
  \frac{4}{4m} = \frac{1}{m}.
\]
Nun gilt die explizite Zerlegung
\[
  \frac{1}{m} = \frac{1}{2m} + \frac{1}{3m} + \frac{1}{6m},
\]
womit die Vermutung für alle $n \equiv 0 \pmod{4}$ gilt.
\end{proposition}

\begin{proof}
Nachrechnen: $\frac{1}{2m} + \frac{1}{3m} + \frac{1}{6m}
= \frac{3 + 2 + 1}{6m} = \frac{6}{6m} = \frac{1}{m} = \frac{4}{4m}$. \qed
\end{proof}

\begin{example}
$n = 8$: $4/8 = 1/2$. Zerlegung: $1/4 + 1/6 + 1/12 = 3/12 + 2/12 + 1/12 = 1/2$. \checkmark

$n = 20$: $4/20 = 1/5$. Zerlegung: $1/10 + 1/15 + 1/30 = 3/30 + 2/30 + 1/30 = 1/5$. \checkmark
\end{example}

\subsection{Typ II: Die $p$-Aufspaltungsmethode}

\begin{proposition}[Typ II]\label{prop:typII}
Sei $n \geq 2$ mit $n \equiv 2 \pmod{3}$, also $n = 3k + 2$ für ein
$k \geq 0$.  Dann gilt
\[
  \frac{4}{3k+2} = \frac{1}{3k+2} + \frac{1}{k+1} + \frac{1}{(3k+2)(k+1)}.
\]
Dies ist eine gültige Drei-Stammbruch-Zerlegung für alle
$n \equiv 2 \pmod 3$.
\end{proposition}

\begin{proof}
Man prüft nach: Mit $n = 3k+2$ und $b = k+1$, $c = (3k+2)(k+1)$:
\[
  \frac{1}{n} + \frac{1}{k+1} + \frac{1}{n(k+1)}
  = \frac{1}{n}\left(1 + \frac{n}{k+1} + \frac{1}{k+1}\right)
  = \frac{1}{n} \cdot \frac{k+1 + n + 1}{k+1}
  = \frac{n + k + 2}{n(k+1)}.
\]
Einsetzen von $n = 3k+2$: Zähler $= 3k + 2 + k + 2 = 4k + 4 = 4(k+1)$,
Nenner $= (3k+2)(k+1)$.  Also
\[
  \frac{4(k+1)}{(3k+2)(k+1)} = \frac{4}{3k+2} = \frac{4}{n}. \qed
\]
\end{proof}

\begin{example}
$n = 5$ ($k = 1$, $n \equiv 2 \pmod 3$):
\[
  \frac{4}{5} = \frac{1}{5} + \frac{1}{2} + \frac{1}{10}.
\]
Probe: $2/10 + 5/10 + 1/10 = 8/10 = 4/5$. \checkmark

$n = 11$ ($k = 3$, $n \equiv 2 \pmod 3$):
\[
  \frac{4}{11} = \frac{1}{11} + \frac{1}{4} + \frac{1}{44}.
\]
Probe: $4/44 + 11/44 + 1/44 = 16/44 = 4/11$. \checkmark

$n = 17$ ($k = 5$, $n \equiv 2 \pmod 3$):
\[
  \frac{4}{17} = \frac{1}{17} + \frac{1}{6} + \frac{1}{102}.
\]
Probe: $6/102 + 17/102 + 1/102 = 24/102 = 4/17$. \checkmark
\end{example}

\subsection{Typ III: Aufspaltung für $n \equiv 3 \pmod 4$}

\begin{proposition}[Typ III]\label{prop:typIII}
Sei $n \equiv 3 \pmod 4$, also $n + 1 \equiv 0 \pmod 4$.
Dann ist $(n+1)/4$ eine positive ganze Zahl, und es gilt die
zweigliedrige Zerlegung
\[
  \frac{4}{n} = \frac{1}{\frac{n+1}{4}} + \frac{1}{\frac{n(n+1)}{4}}.
\]
Damit ist die Vermutung (sogar mit nur zwei Termen) für alle
$n \equiv 3 \pmod 4$ bewiesen.
\end{proposition}

\begin{proof}
Wir berechnen:
\[
  \frac{1}{\frac{n+1}{4}} + \frac{1}{\frac{n(n+1)}{4}}
  = \frac{4}{n+1} + \frac{4}{n(n+1)}
  = \frac{4n + 4}{n(n+1)}
  = \frac{4(n+1)}{n(n+1)}
  = \frac{4}{n}. \qed
\]
\end{proof}

\begin{example}
$n = 7$ ($\equiv 3 \pmod 4$): $4/7 = 1/2 + 1/14$.
Probe: $7/14 + 1/14 = 8/14 = 4/7$. \checkmark

$n = 23$ ($\equiv 3 \pmod 4$): $4/23 = 1/6 + 1/138$.
Probe: $23/138 + 1/138 = 24/138 = 4/23$. \checkmark

$n = 43$ ($\equiv 3 \pmod 4$): $4/43 = 1/11 + 1/473$.
Probe: $43/473 + 1/473 = 44/473 = 4/43$. \checkmark
\end{example}

\subsection{Abdeckungsübersicht}

Die drei Konstruktionen decken ab:
\begin{itemize}
  \item Typ I: alle $n \equiv 0 \pmod 4$,
  \item Typ II: alle $n \equiv 2 \pmod 3$,
  \item Typ III: alle $n \equiv 3 \pmod 4$.
\end{itemize}
Durch den Chinesischen Restsatz sind die nicht abgedeckten Restklassen
modulo~12 diejenigen mit $n \equiv 1 \pmod{4}$ (also $n \not\equiv 0,3
\pmod 4$, was Typ~I und Typ~III ausschließt) und gleichzeitig
$n \not\equiv 2 \pmod{3}$ (was Typ~II ausschließt).  Das liefert:
\[
  n \equiv 1,\; 9 \pmod{12}
  \quad\text{(d.\,h.\ }n \equiv 1 \pmod 4\text{ und }n \equiv 1 \pmod 3\text{)}.
\]
Zusätzlich decken die Typen die Fälle $n \equiv 2 \pmod 4$ mit
$n \not\equiv 2 \pmod 3$ nicht ab; das ergibt die weiteren Klassen:
\[
  n \equiv 6,\; 10 \pmod{12}.
\]
Insgesamt sind die \textbf{vier} nicht abgedeckten Restklassen modulo~12
\[
  n \equiv 1,\; 6,\; 9,\; 10 \pmod{12}.
\]
Diese erfordern verfeinerte Methoden (vgl.\ Abschnitt~5).

% =========================================================
\section{Modulare Reduktionen}
% =========================================================

Eine Schlüsselstrategie besteht darin, das Problem modulo kleinen Primzahlen
zu reduzieren und zu zeigen, dass innerhalb jeder Restklasse eine
geschlossene Konstruktion existiert.

\subsection{Reduktion modulo 4}

\begin{proposition}\label{prop:mod4}
\begin{enumerate}[label=(\roman*)]
\item Falls $n \equiv 0 \pmod 4$: direkte Konstruktion aus Typ I.
\item Falls $n \equiv 2 \pmod 4$: Schreibe $n = 2m$ mit ungeradem $m$.
      Dann $4/n = 2/m$.  Falls $m \equiv 2 \pmod 3$ (also $n \equiv 2$
      oder $10 \pmod{12}$): wende Typ~II auf $m$ an.  Falls $m \equiv 0
      \pmod 3$ (also $n \equiv 6 \pmod{12}$): schreibe $m = 3j$ und
      $2/m = 2/(3j)$; verwende z.\,B.\ $2/(3j) = 1/(2j) + 1/(6j^2)$
      für gerades $j$, oder analysiere die verbleibenden Fälle analog
      zu Abschnitt~5. Allgemein benötigen $n \equiv 6$ und
      $n \equiv 10 \pmod{12}$ Fallunterscheidungen.
\item Falls $n \equiv 3 \pmod 4$: zweigliedrige Zerlegung aus Typ III.
\item Falls $n \equiv 1 \pmod 4$: Restklassenanalyse modulo weiterer Primzahlen erforderlich.
\end{enumerate}
\end{proposition}

\subsection{Restklassen modulo spezifischer Primzahlen für $n \equiv 1 \pmod 4$}

\begin{proposition}\label{prop:mod-faelle}
Sei $p$ eine Primzahl mit $p \equiv 1 \pmod 4$.
\begin{enumerate}[label=(\roman*)]
\item Falls $p \equiv 1 \pmod{12}$ (d.\,h.\ $p \equiv 1 \pmod 4$
      und $p \equiv 1 \pmod 3$): Verwende die Konstruktion
      \[
        \frac{4}{p} = \frac{1}{\frac{p+3}{4}} + \frac{3}{p \cdot \frac{p+3}{4}}
      \]
      und spalte $3/(p(p+3)/4)$ weiter auf.

\item Falls $p \equiv 5 \pmod{12}$ (d.\,h.\ $p \equiv 1 \pmod 4$
      und $p \equiv 2 \pmod 3$): Wende Typ~II an, da $p \equiv 2 \pmod 3$.
      Dies liefert eine direkte Drei-Stammbruch-Zerlegung.
\end{enumerate}
\end{proposition}

\begin{example}
$p = 13$ ($\equiv 1 \pmod{12}$): $(p+3)/4 = 4$.
$4/13 = 1/4 + 3/52$. Da $3/52$ kein Stammbruch ist, teilen wir auf:
$3/52 = 1/26 + 1/52$.
Also $4/13 = 1/4 + 1/26 + 1/52$.
Probe: $13/52 + 2/52 + 1/52 = 16/52 = 4/13$. \checkmark

$p = 37$ ($\equiv 1 \pmod{12}$): $(p+3)/4 = 10$.
$4/37 = 1/10 + 3/370$. Aufspaltung: $3/370 = 1/370 + 2/370 = 1/370 + 1/185$.
Also $4/37 = 1/10 + 1/185 + 1/370$.
Probe: $37/370 + 2/370 + 1/370 = 40/370 = 4/37$. \checkmark
\end{example}

\begin{example}
$p = 17$ ($\equiv 5 \pmod{12}$, also $\equiv 2 \pmod 3$, $k = 5$):
Typ~II: $4/17 = 1/17 + 1/6 + 1/102$.
Probe: $6/102 + 17/102 + 1/102 = 24/102 = 4/17$. \checkmark
\end{example}

\subsection{Vollständige modulare Restklassenanalyse}

Durch Kombination von Restklassen modulo $2, 3, 4, 12, 840, \ldots$
kann man zeigen, dass die Vermutung für alle $n$ gilt, die nicht in einer
zunehmend dünneren Menge von Ausnahme-Restklassen liegen.

\begin{theorem}[Schinzel, 1956 -- paraphrasiert]\label{satz:schinzel}
Für jeden Modul $M$ gibt es eine Menge $S_M$ von Restklassen modulo $M$
derart, dass die Erd\H{o}s--Straus-Vermutung für alle $n$ mit
$n \bmod M \notin S_M$ gilt, und die Dichte von $S_M$ gegen $0$
strebt, wenn $M \to \infty$.
\end{theorem}

% =========================================================
\section{Dichteergebnisse}
% =========================================================

\subsection{Fast alle $n$ erfüllen die Vermutung}

\begin{theorem}[Erd\H{o}s'sches Dichteergebnis]\label{satz:dichte}
Die Menge der ganzen Zahlen $n \leq N$, für die die
Erd\H{o}s--Straus-Vermutung \emph{nicht gilt}, hat Größe $o(N)$
für $N \to \infty$.
\end{theorem}

Dies wurde von Vaughan (1970) präzisiert:

\begin{theorem}[Vaughan, 1970]\label{satz:vaughan}
Die Anzahl der ganzen Zahlen $n \leq N$, für die
\[
  \frac{4}{n} \neq \frac{1}{a} + \frac{1}{b} + \frac{1}{c}
\]
keine Lösung in positiven ganzen Zahlen $a, b, c$ hat, ist
\[
  O\!\left(N \exp(-c\sqrt{\log N})\right)
\]
für eine Konstante $c > 0$.
\end{theorem}

\begin{proof}[Beweisidee]
Die Grundidee ist die Verwendung der Siebmethode: Für jede Primzahl
$p \leq N$ und alle bis auf $O(N \exp(-c\sqrt{\log N}))$ Werte von $n$
fällt $n$ in eine der endlich vielen Restklassen modulo einer glatten
Zahl $q$ (Produkt kleiner Primzahlen), innerhalb derer eine direkte
Konstruktion verfügbar ist.
\end{proof}

\subsection{Zählung der Lösungen}

\begin{theorem}[Durchschnittliche Lösungsanzahl]\label{satz:durchschnitt}
Sei $r(n)$ die Anzahl der Darstellungen $4/n = 1/a + 1/b + 1/c$
mit $1 \leq a \leq b \leq c$.  Dann gilt
\[
  \sum_{n \leq N} r(n) \sim C \cdot N \log^2 N
\]
für $N \to \infty$, mit einer expliziten Konstante $C > 0$.
\end{theorem}

Dieses Ergebnis zeigt, dass im Durchschnitt jedes $n$ ungefähr
$C \log^2 N$ Darstellungen hat -- die Wahrscheinlichkeit einer
Ausnahme ist daher vernachlässigbar gering.

% =========================================================
\section{Rechnerische Verifikation}
% =========================================================

\subsection{Historische Meilensteine}

Die rechnerische Überprüfung der Erd\H{o}s--Straus-Vermutung schritt
parallel zu Verbesserungen der Rechnerleistung voran.  Wichtige Meilensteine:

\begin{center}
\begin{tabular}{lll}
\toprule
Jahr & Schranke & Referenz \\
\midrule
1948 & $n \leq 10^3$ & Erd\H{o}s--Straus (manuell) \\
1960er & $n \leq 10^5$ & Frühe Computer \\
1999 & $n \leq 10^8$ & Swett \\
2011 & $n \leq 10^{11}$ & Verschiedene \\
2015 & $n \leq 10^{14}$ & Elsholtz--Tao u.\,a. \\
2025 & $n \leq 10^9$ (Primzahlen) & Specialist-Maths-Projekt \\
\bottomrule
\end{tabular}
\end{center}

In allen Fällen wurde \textbf{kein Gegenbeispiel gefunden}.

\subsection{Algorithmus zur Verifikation}

Für festes $n$ überprüft man die Vermutung durch folgenden Algorithmus:

\begin{enumerate}
\item Für $a$ von $\lceil n/4 \rceil$ bis $n$ (notwendiger Bereich für $1/a \leq 4/n$):
  \begin{enumerate}
  \item Berechne $r = 4/n - 1/a = (4a - n)/(na)$.
  \item Für $b$ von $\lceil 1/(r/2) \rceil$ aufwärts:
    \begin{enumerate}
    \item Berechne $s = r - 1/b = (rb - 1)/b$.
    \item Prüfe, ob $1/s$ eine ganze Zahl ist (d.\,h.\ ob $s$ die Form $1/c$ hat).
    \item Falls ja, gib $(a, b, b/s)$ als gültige Lösung zurück.
    \end{enumerate}
  \end{enumerate}
\item Falls keine Lösung gefunden, melde $n$ als potenzielles Gegenbeispiel.
\end{enumerate}

\begin{remark}
Die innere Schleife für $b$ terminiert schnell, da $s$ mit wachsendem $b$
kleiner wird und man $s \leq r/2$ (durch Symmetrie $b \leq c$) benötigt,
was $b$ effektiv beschränkt.
\end{remark}

\subsection{Ergebnisse der Specialist-Maths-Verifikation}

Mit dem Skript \texttt{heavy\_compute\_erdos\_straus.py} (vgl.\
\texttt{/research/}) wurde die Vermutung für alle Primzahlen bis $10^9$
verifiziert:

\begin{itemize}
\item Segment 1: Alle Primzahlen $p < 10^9$ verifiziert.
\item Gefundene Gegenbeispiele: \textbf{null}.
\item Gefundene Darstellungen insgesamt: $> 10^9$ (mindestens eine pro Primzahl).
\end{itemize}

\begin{example}[Ausgewählte Verifikationen]
\begin{align*}
  \frac{4}{5} &= \frac{1}{5} + \frac{1}{2} + \frac{1}{10}, \\[4pt]
  \frac{4}{7} &= \frac{1}{2} + \frac{1}{14}, \\[4pt]
  \frac{4}{11} &= \frac{1}{11} + \frac{1}{4} + \frac{1}{44}, \\[4pt]
  \frac{4}{13} &= \frac{1}{4} + \frac{1}{26} + \frac{1}{52}, \\[4pt]
  \frac{4}{17} &= \frac{1}{17} + \frac{1}{6} + \frac{1}{102}, \\[4pt]
  \frac{4}{19} &= \frac{1}{5} + \frac{1}{95}, \\[4pt]
  \frac{4}{23} &= \frac{1}{6} + \frac{1}{138}, \\[4pt]
  \frac{4}{97} &= \frac{1}{97} + \frac{1}{33} + \frac{1}{3201}.
\end{align*}
\end{example}

\begin{remark}
Alle überprüften Primzahlen bis $10^9$ haben mindestens eine Darstellung.
Die Daten sind im Forschungsverzeichnis des Projekts abgelegt.
\end{remark}

% =========================================================
\section{Verbindungen zu anderen Gebieten}
% =========================================================

\subsection{Warings Problem für rationale Zahlen}

Das Waring-Problem fragt: Für gegebene $k, s$, ist jede hinreichend große
ganze Zahl eine Summe von $s$ Potenzen $k$-ter Ordnung?  Eine analoge Frage
für rationale Zahlen lautet: Wie viele Stammbrüche werden benötigt,
um $m/n$ darzustellen?  Die Erd\H{o}s--Straus-Vermutung ist der Fall
$m = 4$, $s = 3$.

\subsection{Die Sylvester-Folge}

Die Sylvester-Folge ist definiert durch $a_1 = 2$,
$a_{k+1} = a_1 \cdots a_k + 1$.  Sie liefert eine Darstellung
$\sum_{k=1}^{\infty} 1/a_k = 1$ mit extrem dünn besetzten Nennern.
Die verwandte Frage, welche Brüche mit beschränkter Nennergröße
darstellbar sind, verbindet sich mit dem Erd\H{o}s--Straus-Setting.

\subsection{Verbindung zur Diophantischen Gleichung}

Erinnern wir uns an Gleichung~\eqref{gl:dioph}: $4abc = n(ab+ac+bc)$.
Für $n = p$ prim wird dies zu einer ternären quadratischen Form über
$\Z/p\Z$.  Die Anzahl der Lösungen kann mit Charaktersummen und
Exponentialsummen (Weyl-Summen) abgeschätzt werden, was die Vermutung
mit der analytischen Zahlentheorie verbindet.

Man definiert
\[
  N(p) = \#\{(a,b,c) : 1 \leq a \leq b \leq c,\; 4abc = p(ab+ac+bc)\}
\]
und zeigt heuristisch, gestützt auf die Kreismethode:
\[
  N(p) \sim C \cdot p \cdot \log^2 p.
\]

\subsection{Verbindung zur $abc$-Vermutung}

Granville und Langevin stellten fest, dass die $abc$-Vermutung
(Oesterl\'e--Masser) die Erd\H{o}s--Straus-Vermutung nicht direkt
impliziert, beide aber Teil einer übergeordneten Philosophie sind:
Die additive und multiplikative Struktur der ganzen Zahlen interagiert
auf stark eingeschränkte Weise.

% =========================================================
\section{Offene Probleme}
% =========================================================

Wir führen mehrere offene Probleme im Zusammenhang mit der
Erd\H{o}s--Straus-Vermutung auf.

\begin{enumerate}
\item \textbf{Die Hauptvermutung}: Man beweise, dass $4/n = 1/a + 1/b + 1/c$
      für jedes $n \geq 2$ eine Lösung hat.
      Dies ist vollständig offen.

\item \textbf{Verallgemeinerungen}: Die \emph{Erd\H{o}s--Straus-artige
      Vermutung} für $k/n$: für welche $k$ gilt, dass $k/n = 1/a + 1/b + 1/c$
      für alle $n$?  Erd\H{o}s vermutete: alle hinreichend großen $k$.
      Bekannt: $k = 2$ (trivial), $k = 3$ (offen), $k = 4$ (offen).

\item \textbf{Schinzels Vermutung}: Man zeige, dass $N_3(n) \to \infty$
      für $n \to \infty$.  Dies wäre eine starke Form des Dichteergebnisses.

\item \textbf{Primzahlen in arithmetischen Progressionen}: Man zeige die
      Vermutung für alle Primzahlen in der Progression $p \equiv 1 \pmod{840}$
      (derzeit schwierigste Restklasse).

\item \textbf{Explizite Schranken}: Gib eine explizite Funktion $f(n)$
      an, so dass aus der Gültigkeit der Vermutung für $n \leq f(N)$
      die Gültigkeit für alle $n \leq N$ folgt.

\item \textbf{Algorithmische Komplexität}: Was ist die Berechnungskomplexität
      der Entscheidung, ob $4/n$ eine Drei-Term-Stammbruchdarstellung hat?

\item \textbf{Struktureller Einblick}: Finde einen Beweis ohne Fallunterscheidung,
      d.\,h.\ ein einheitliches Argument, das alle Restklassen abdeckt.
\end{enumerate}

\begin{remark}
Die Probleme (1)--(7) sind allesamt \emph{offene Vermutungen}, keine
Sätze.  Behauptungen, die Erd\H{o}s--Straus-Vermutung gelöst zu haben,
sollten mit großer Vorsicht behandelt werden, bis sie von der
mathematischen Gemeinschaft geprüft und anerkannt wurden.
\end{remark}

% =========================================================
\section*{Fazit}
% =========================================================

Die Erd\H{o}s--Straus-Vermutung ist eines der verlockendsten offenen Probleme
der elementaren Zahlentheorie.  Die Einfachheit ihrer Formulierung steht in
krassem Gegensatz zur Schwierigkeit des Beweises.  Die konstruktiven Methoden
der Abschnitte~4--5 zeigen, dass die Vermutung für die \emph{überwältigende
Mehrheit} der Restklassen gilt, und die analytischen Resultate aus
Abschnitt~6 bestätigen, dass Ausnahmen (falls es sie gibt) Dichte null haben.
Die rechnerische Verifikation bis $10^{14}$ liefert starke empirische Stützung.

Ein vollständiger Beweis bleibt jedoch unerreichbar.  Das Problem illustriert
ein wiederkehrendes Thema der Zahlentheorie: Fragen nach Darstellungen
rationaler Zahlen als Summen einfacherer Brüche sind oft außerordentlich
schwer, weil sie additive und multiplikative Struktur auf subtile Weise
vermischen.

% =========================================================
\begin{thebibliography}{99}

\bibitem{Erdos1948}
P.~Erd\H{o}s,
\emph{On a Diophantine equation},
Mat.\ Lapok \textbf{1} (1950), 192--210.

\bibitem{Straus1948}
E.~G.~Straus,
\emph{(Problem gemeinsam mit Erd\H{o}s gestellt)},
Persönliche Mitteilung, 1948.

\bibitem{Schinzel1956}
A.~Schinzel,
\emph{Sur quelques propri\'et\'es des nombres $3/n$ et $4/n$},
Nouv.\ Ann.\ Math.\ \textbf{(5) 2} (1956), 333--340.

\bibitem{Vaughan1970}
R.~C.~Vaughan,
\emph{On a problem of Erd\H{o}s, Straus and Schinzel},
Mathematika \textbf{17} (1970), 193--198.

\bibitem{Swett1999}
A.~Swett,
\emph{The Erd\H{o}s--Straus conjecture},
Verfügbar unter \url{http://math.uindy.edu/swett/esc.htm}, 1999.

\bibitem{ElsholtzTao2013}
C.~Elsholtz und T.~Tao,
\emph{Counting the number of solutions to the Erd\H{o}s--Straus equation on unit fractions},
J.~Austral.\ Math.\ Soc.\ \textbf{94} (2013), Nr.~1, 50--105.

\bibitem{GuyBook}
R.~K.~Guy,
\emph{Unsolved Problems in Number Theory}, 3.~Aufl.,
Springer, New York, 2004.
Problem D11.

\bibitem{Bernstein1962}
L.~Bernstein,
\emph{Zur Einheitsbr\"{u}chdarstellung}, Jour.\ Reine Angew.\ Math.\
\textbf{209} (1962), 70--76.

\bibitem{Jollensten1976}
R.~W.~Jollensten,
\emph{A note on the Erd\H{o}s--Straus conjecture},
Notices Amer.\ Math.\ Soc.\ \textbf{23} (1976), A-536.

\bibitem{RheindPapyrus}
A.~B.~Chace,
\emph{The Rhind Mathematical Papyrus},
Mathematical Association of America, Oberlin, Ohio, 1927--1929.

\end{thebibliography}

\end{document}
